\section*{第 \ref{ch:g11:dissolution} 章——溶解离子固体和分子固体}

\begin{solution}{exo:g11:dissolution:1}
\ce{NaCl(s) -> Na+(aq) + Cl-(aq)};
\ce{CaCl2(s) -> Ca^{2+}(aq) + 2Cl-(aq)};
\ce{Na2SO4(s) -> 2Na+(aq) + SO4^{2-}(aq)};
\ce{AlCl3(s) -> Al^{3+}(aq) + 3Cl-(aq)}.
\end{solution}

\begin{solution}{exo:g11:dissolution:2}
$[\ce{Na+}] = 2 \times 0.050 = \qty{0.10}{mol/L}$;
$[\ce{SO4^{2-}}] = \qty{0.050}{mol/L}$.
\end{solution}

\begin{solution}{exo:g11:dissolution:3}
水合是指离子被与其相吸引的水分子包围。
氧端（$\delta^-$）指向阳离子，氢端
（$\delta^+$）指向阴离子：异种电荷相互吸引。
\end{solution}

\begin{solution}{exo:g11:dissolution:4}
$M(\ce{AlCl3}) = 27.0 + 3 \times 35.5 = \qty{133.5}{g/mol}$;
$n = 2.67 / 133.5 = \qty{0.0200}{mol}$; $c = 0.0200 / 0.2000 =
\qty{0.100}{mol/L}$; $[\ce{Al^{3+}}] = \qty{0.100}{mol/L}$,
$[\ce{Cl-}] = \qty{0.300}{mol/L}$.
\end{solution}

\begin{solution}{exo:g11:dissolution:5}
头部 \ce{-COO-} 为离子性基团，受水吸引，属亲水的。尾部
由 17 个碳原子组成，仅含 \ce{C-C} 键和近乎非极性的 \ce{C-H}
键，属非极性部分，故为疏水的。
\end{solution}

\begin{solution}{exo:g11:dissolution:6}
盐：水（离子固体，极性溶剂）。蜡：环己烷（非极性
链）。糖：水（其 \ce{O-H} 基团与
水形成氢键）。碘：环己烷（非极性分子）。
\end{solution}

\begin{solution}{exo:g11:dissolution:7}
乙醇的 \ce{O-H} 基团与水分子形成氢键，
所形成的氢键替代其所断开的氢键：乙醇能融入水中。己烷为非极性，
不能形成此类氢键；水分子借彼此间的氢键相连，
因而不允许其进入。
\end{solution}

\begin{solution}{exo:g11:dissolution:8}
疏水尾部避开水，溶于油脂之中，位于
胶束中心；带电的头部受水吸引，留在外侧。所有
胶束表面都带负电，同种电荷相互排斥，
因而彼此保持分开。
\end{solution}

\begin{solution}{exo:g11:dissolution:9}
环己烷与水不互溶，因而形成一独立液层，
可以分液放出；非极性的香料在其中溶解良好。
乙醇与水互溶，不会形成第二液层，
因而无法将其分离。
\end{solution}

\begin{solution}{exo:g11:dissolution:10}
环己烷的密度小于水。下层的水相
经活塞放出。碘位于上层的环己烷层中，该层
留在分液漏斗中。
\end{solution}

\begin{solution}{exo:g11:dissolution:11}
无溶解：硫酸铜为离子化合物，不溶于非极性的
环己烷。水仍呈蓝色，环己烷层无色。
\end{solution}

\begin{solution}{exo:g11:dissolution:12}
$[\ce{Cl-}] = 2c$，故 $c = \qty{0.150}{mol/L}$；
$n = 0.150 \times 0.2500 = \qty{0.0375}{mol}$;
$m = 0.0375 \times 111.1 = \qty{4.17}{g}$.
\end{solution}

\begin{solution}{exo:g11:dissolution:13}
钙离子与镁离子使硬脂酸根离子以浮渣形式沉淀：
肥皂在形成胶束前即被耗尽，起泡效果很差。
仅就钙而言：\qty{0.010}{mol} 的 \ce{Ca^{2+}} 需要
\qty{0.020}{mol} 的硬脂酸根，即 $0.020 \times 306.0 =
\qty{6.1}{g}$ 的硬脂酸钠（每升）。
\end{solution}

\begin{solution}{exo:g11:dissolution:14}
体积 \qty{0.2000}{L}。\ce{Na+}：$0.020 / 0.2000 = \qty{0.10}{mol/L}$；
\ce{Ca^{2+}}: $0.010 / 0.2000 = \qty{0.050}{mol/L}$; \ce{Cl-}:
$(0.020 + 0.020) / 0.2000 = \qty{0.20}{mol/L}$。正电荷：
$0.10 + 2 \times 0.050 = 0.20$；负电荷：$0.20$。呈电中性。
\end{solution}

\begin{solution}{exo:g11:dissolution:15}
一升庚烷的质量为 \qty{0.68}{kg}，可溶解
$17 \times 0.68 = \qty{12}{g}$ 的碘，
约为一升水的 40 倍。混合振荡后，几乎全部碘转入
庚烷中，萃取效率很高。
\end{solution}

\begin{solution}{pb:g11:dissolution:1}
\textbf{1.} 水流经的岩石（石灰石、
白云石、石膏），其中少量溶解。

\textbf{2.} $[\ce{Ca^{2+}}] = 0.120 / 40.1 = \qty{2.99e-3}{mol/L}$;
$[\ce{Mg^{2+}}] = 0.024 / 24.3 = \qty{9.9e-4}{mol/L}$.

\textbf{3.} \ce{Ca(HCO3)2 -> Ca^{2+}(aq) + 2HCO3-(aq)}；故
$[\ce{HCO3-}] = 2 \times \num{2.99e-3} = \qty{5.99e-3}{mol/L}$.

\textbf{4.} $\num{2.99e-3} \times 150 = \qty{0.449}{mol}$.

\textbf{5.} $18 \times 12.0 + 35 \times 1.0 + 2 \times 16.0 + 23.0 =
\qty{306.0}{g/mol}$.

\textbf{6.} \ce{C17H35COONa(s) -> C17H35COO-(aq) + Na+(aq)}。

\textbf{7.} 头部 \ce{-COO-}：亲水的；尾部 \ce{C17H35-}：
疏水的。它兼具两种部分，是两亲的。

\textbf{8.} 一个微小的球体，由肥皂离子构成，尾部朝内，头部朝外。
尾部溶解于油脂中，油脂被分割成液滴，
液滴被包裹在胶束中；带电的表面使液滴在水中相互分开，
从而被水冲走。

\textbf{9.} 肥皂必须形成离子和胶束，这需要头部周围有水；
油脂则被冲洗用的水带走。在
油中，油脂只会留在皮肤上。

\textbf{10.} \ce{2C17H35COO- + Ca^{2+} -> Ca(C17H35COO)2}：
电荷：左边 $2 \times (-1) + 2 = 0$，右边 $0$。

\textbf{11.} \ce{C17H35COO-}：剩余 $- 2x$；\ce{Ca^{2+}}：
$\num{2.99e-3} - x$；硬脂酸钙：$x$。钙离子为
限量离子：$x_{\max} = \qty{2.99e-3}{mol}$。

\textbf{12.} $2 \times \num{2.99e-3} = \qty{5.99e-3}{mol}$ 每升。

\textbf{13.} $M = 36 \times 12.0 + 70 \times 1.0 + 4 \times 16.0 + 40.1 =
\qty{606.1}{g/mol}$；$\num{2.99e-3} \times 606.1 = \qty{1.81}{g}$ 的浮渣
每升。

\textbf{14.} $2 \times \num{9.9e-4} = \qty{1.98e-3}{mol}$ 每升。

\textbf{15.} $2 \times 0.449 = \qty{0.898}{mol}$。

\textbf{16.} $0.898 \times 306.0 = \qty{275}{g}$。

\textbf{17.} 镁：$\num{9.9e-4} \times 150 = \qty{0.148}{mol}$，
需要 \qty{0.296}{mol} 硬脂酸根，\qty{91}{g} 肥皂；合计
约 $275 + 91 = \qty{366}{g}$。

\textbf{18.} $275 / 100 \approx 2.7$ 巴。

\textbf{19.} 一浴水中的钙约浪费 \qty{0.27}{kg} 肥皂
变成浮渣。
\end{solution}
